\documentclass{article}
\title{ECE 478 -- Midterm 3 Practice}
\usepackage[margin=1in]{geometry}
\usepackage{graphicx}
\usepackage{pdfpages}
\usepackage{fancyhdr}
\usepackage{enumitem}
\usepackage{pdfpages}
\usepackage{amsmath}
\usepackage{amssymb}
\usepackage{float}
\pagestyle{fancy}
\fancyhf{}
\fancyhead[L]{Gianni Avilan-Losee}
\fancyhead[C]{ECE 478}
\fancyhead[R]{Midterm 2 Practice}
\fancyfoot[C]{\thepage}
\usepackage{microtype}
\usepackage{textcomp, gensymb}
\author{Gianni Avilan-Losee}
\begin{document}
\setcounter{section}{5}
\section{Chapter 6}
\subsection{Review}
Wire geometry, resistance, capacitance, inductance. 
\(\pi\)-model for finding Elmore delay. Use sheet resistance to eliminate detail and find resistance per unit length, using the width. Scaling affects RC wire delay. Crosstalk delay, Miller effect (\textbf{note:} definitely look further into Miller effect.) Capacitor voltage divider. Repeaters! Best length of wire between inverters, best repeater size. 
\subsection{Questions}
Pi-model question, resistivity question (sheet), multiple-choice scaling/RC question, Miller effect, maybe repeater.
\section{Chapter 7}
\subsection{Review}
Process variation can be rolled into one variation parameter, \(K_{\mathrm{process}}\), so we don't need to labor over the details. Note that \(K_{\mathrm{load}}\) is given in \(\mathrm{\dfrac{ns}{pF}}\), so find the wire + gate capacitance of the load and multiply by \(K_{\mathrm{load}}\). Dennard scaling is the simplest assumption, where all dimensions scale in the same way to ensure a constant electric field. 
\subsection{Questions}
TSMC cell library derating factors, effects of voltage, temperature, process (all possible manufacturing factors rolled into one). Might include loaded gates. Definitely a Dennard scaling question.
\section{Chapter 9}
\subsection{Review}
Inner/outer input issues (\textbf{review}). Why does this happen? Thinking about a series nMOS network, say \(B\) and \(C\) inputs are already 1 (the only condition in which a rising edge on \(A\) could case a falling \(Y\) output). Generalize this to a NOR gate? Same result? No, but we might have a similar situation on the rising edge.

Say \(B\) transition causes \(Y\) to rise, then \(A\) must already be \(0\), and the \(A\) capacitor pre-charged. In the alternate case, where \(A\) transistions \textit{last}, both capacitors will need to charge. Asymmetric gates are complex to analyze, but the basic rationale makes sense. Skewed gates, \textbf{review}. Gets away from matching pull-up and pull-down strengths (why would you do this?) Ratioed circuits use an always-on pMOS transistor (basically a resistor). Key feature is the ratio of pull-up current to pull-down current (\textbf{review}). F e e t. Monotonic inputs, you can repeat these to make them work, but you can't cascade \textit{dynamic} gates.
\subsection{Questions}
Three-input multiple choice, NAND/NOR. Fill in transistor sizing for HI-skew and LO-skew (multiple-choice?). Probably something on ratioed and/or f e e t.
\end{document}
